%!TEX root = ../graduate-work.tex
\section{Parallelism Strategies for Neural Network Training}

The parallelism methods examined in this work were developed for training large
language models (LLMs) that do not fit in the memory of a single GPU. Three
principal strategies are described below---Tensor Parallelism, Pipeline
Parallelism, and Fully Sharded Data Parallelism---along with the rationale for
considering them as candidates for adaptation to the verification problem.

\subsection{Scale of Modern Models and the Transformer Architecture}

The backbone of modern LLMs is the Transformer architecture~\cite{Vaswani2017Attention},
built around the self-attention mechanism. Each layer consists of a
Multi-Head Attention block and a Feed-Forward Network.

In the \textbf{Multi-Head Attention} block, query, key, and value matrices are
computed as $Q = XW_Q$, $K = XW_K$, $V = XW_V$, where
$W_Q, W_K, W_V \in \mathbb{R}^{d \times d}$, and the attention output is:
$$\mathrm{Attention}(Q, K, V) = \mathrm{softmax}\!\left(\frac{QK^\top}{\sqrt{d}}\right)V.$$

The \textbf{Feed-Forward Network} is a two-layer fully connected network with
an expansion factor:
$$\mathrm{FFN}(x) = W_2 \cdot \mathrm{GELU}(W_1 x + b_1) + b_2,$$
where $W_1 \in \mathbb{R}^{4d \times d}$ and $W_2 \in \mathbb{R}^{d \times 4d}$.

Each block contains several matrix multiplications with weight matrices of
size $d \times d$ or $d \times 4d$. At $d = 12288$ (GPT-3, 175 billion
parameters~\cite{GPT3Brown2020}), a single matrix $W_1$ occupies
$12288 \times 49152 \times 4 \approx 2.3$~GB in float32. A model with
96 such layers does not fit on any existing accelerator---hence the need
for parallelism strategies.

\subsection{Tensor Parallelism (TP)}

Tensor Parallelism, introduced in Megatron-LM~\cite{MegatronLM2020},
partitions matrix multiplications \textit{within a layer} across multiple GPUs.
The weight matrix is split so that each GPU performs its portion of the
multiplication, and results are synchronised via collective operations.

\subsubsection{Column Parallel Linear}

The weight matrix $W \in \mathbb{R}^{k \times n}$ is partitioned along the
column (output) dimension into $P$ shards:
$$W = [W^{(0)} \mid W^{(1)} \mid \ldots \mid W^{(P-1)}], \quad
  W^{(r)} \in \mathbb{R}^{k \times (n/P)}.$$
GPU $r$ stores only its shard $W^{(r)}$ and computes
$\mathbf{y}^{(r)} = W^{(r)\top} \mathbf{x}$, where the input vector
$\mathbf{x}$ is replicated on all GPUs. After the multiplication, the output
is partitioned along the output dimension.

\subsubsection{Row Parallel Linear}

The weight matrix $W \in \mathbb{R}^{n \times m}$ is partitioned along rows:
$$W = \begin{bmatrix} W^{(0)} \\ W^{(1)} \\ \vdots \\ W^{(P-1)} \end{bmatrix},
  \quad W^{(r)} \in \mathbb{R}^{(n/P) \times m}.$$
GPU $r$ multiplies its portion of the split input $\mathbf{x}^{(r)}$
by its shard: $\mathbf{z}^{(r)} = W^{(r)} \mathbf{x}^{(r)}$.
The full output is reconstructed via \texttt{AllReduce}:
$$\mathbf{y} = \sum_{r=0}^{P-1} \mathbf{z}^{(r)}.$$

\subsubsection{Collective Operations}

A Column--Row Parallel pair forms the basic TP building block.
The forward pass requires only one \texttt{AllReduce} after the Row Parallel
layer; the backward pass requires one \texttt{AllReduce} for gradients before
the Column Parallel layer. On fast inter-GPU interconnects (NVLink, NVSwitch),
synchronisation takes a negligible fraction of compute
time~\cite{MegatronLM2020}. Without direct GPU-to-GPU connections (e.g.\ over
PCIe), \texttt{AllReduce} becomes the bottleneck and the gains from partitioning
disappear---hence the practical limit: TP is sensibly applied within a single
node, typically up to 8~GPUs.

\subsection{Pipeline Parallelism (PP)}

Pipeline Parallelism~\cite{GPipePaper} partitions the model differently---not
within layers but between them. Each GPU receives a \textit{stage}: several
consecutive layers. Data flows through stages as in a pipeline, which creates
the \textit{pipeline bubble} problem: while a GPU processes a new batch, the
previous one is already idle; the same occurs during the backward pass.

The standard remedy is to split the input batch into $M$ micro-batches and
process them without waiting. The bubble fraction is then:
$$r_{\text{bubble}} = \frac{P - 1}{M + P - 1},$$
which approaches zero as $M \gg P$, though at the cost of storing intermediate
activations for all $M$ micro-batches simultaneously. Only activations are
transmitted between stages (point-to-point send/recv), making PP more tolerant
of slow inter-GPU links than TP.

\subsection{Fully Sharded Data Parallelism (FSDP)}

Classic data parallelism replicates the entire model on each GPU, partitioning
only data and averaging gradients via \texttt{AllReduce} after each step.
The drawback is that each GPU holds a full copy of the model, gradients, and
optimiser states, consuming large amounts of redundant memory.

Fully Sharded Data Parallelism (FSDP), based on the ZeRO~\cite{ZeROPaper}
framework, shards all three components: each GPU stores only $1/P$ of each
tensor. Before computing a particular layer, an \texttt{AllGather} assembles
the full weights; after the backward pass, FSDP aggregates gradients via
\texttt{ReduceScatter} and applies them independently on each GPU to its shards.
The method is especially effective for training, where optimiser states
(Adam stores first and second moments) consume triple the memory of the weights.

FSDP became the primary adaptation algorithm in this work: during verification
there are no optimiser states, so \texttt{AllGather}/free of shards can be
embedded in the bound-propagation procedure without altering the backward-pass
logic---details in Chapter~\ref{sec:parallel-verifiers}.
